% SEGMENT OLP-0048-S01
% Part: sets-functions-relations
% Chapter: arithmetization
% Section: cauchy

\documentclass[../../../include/open-logic-section]{subfiles}

\begin{document}

\olfileid{sfr}{arith}{cauchy}

\olsection{இணைப்பு: கோஷி தொடர்களாக மெய்யெண்கள்}

\olref[cuts]{sec} இல், மெய்யெண்களை டெடெகிண்ட் வெட்டுகளாகக்
கட்டமைத்தோம். இந்தப் பிரிவில், ஒரு மாற்றுக் கட்டமைப்பை விளக்குகிறோம்.
அது கோஷி (இப்போது நாம் அழைப்பதன்படி) கோஷி தொடருக்குக் கொடுத்த
வரையறையின் அடிப்படையில் அமைகிறது;
ஆனால் இந்த வரையறையைப் பயன்படுத்தி மெய்யெண்களைக் \emph{கட்டமைத்தவர்கள்},
குறிப்பாக வயர்ஸ்ட்ராஸ், ஹைனே, மெரே, காண்டோர் ஆகிய
19ஆம் நூற்றாண்டின் பிற ஆசிரியர்கள்.
(சிறந்த வரலாற்றுக்கு, \citeauthor{OConnorRobertson:RN}
\citeyear{OConnorRobertson:RN} ஐப் பார்க்கவும்.)

% SEGMENT OLP-0048-S02
19ஆம் நூற்றாண்டை அடைவதற்கு முன், சைமன் ஸ்டெவினை (1548--1620)
கருதுவது பயனுள்ளது.
சுருக்கமாக, ஒவ்வொரு மெய்யெண்ணையும் அதன் தசம விரிவாக்கத்தைக் கொண்டு
கருதலாம் என்பதை ஸ்டெவின் உணர்ந்தார்.
எனவே $\sqrt{2}$ போன்ற ஒரு விகிதமுறா எண்ணுக்குக்கூட,
பின்வருமாறு தொடங்கும் ஒரு நேர்த்தியான தசம விரிவாக்கம் உள்ளது:
\[
	1.41421356237\ldots
\]
தசம விரிவாக்கங்களைக் கணக்கோட்பாட்டில் உருவகப்படுத்துவது மிகவும் எளிது:
அவற்றை $d \colon \Nat \to \Nat$ என்ற சார்புகளாக மட்டும் கருதுக;
நாம் கவனிக்கும் $n$ ஆவது தசம இடம் $d(n)$ ஆகும்.
தசமப் புள்ளிக்கு முன் வரும் மெய்யெண்ணின் பகுதியைக் கையாள
சிறிய மாற்றம் தேவைப்படும் (இங்கு அது $1$ மட்டுமே).
எடுத்துக்காட்டாக, $0.999\ldots = 1$ என்பதை உறுதிசெய்ய,
மேலும் ஒரு மாற்றமும் (ஒரு சமானத் தொடர்பும்) தேவைப்படும்.
ஆனால் ஸ்டெவினின் முறையில், கணக்கோட்பாட்டிற்குள் மெய்யெண்களுக்கான
முழுமையாகக் கடுமையான ஒரு கட்டமைப்பை முன்வைப்பது கடினமல்ல.

% SEGMENT OLP-0048-S03
ஸ்டெவின் நமது கவனப்பொருள் அல்ல.
(ஸ்டெவினைப் பற்றி மேலும் அறிய, \citealt{KatzKatz2012} ஐப் பார்க்கவும்.)
ஆனால் நெருங்கிய தொடர்புடைய ஒரு கருத்து இதோ.
$\sqrt{2}$ இன் தசம விரிவாக்கத்தை நேரடியாகக் கருதுவதற்குப் பதிலாக,
மேன்மேலும் துல்லியமாகும் பின்வரும் விரிவாக்கங்களைக் கொண்டு,
$\sqrt{2}$ ஐ மேன்மேலும் துல்லியமாகத் தோராயப்படுத்தும்
விகிதமுறு எண்களின் ஒரு \emph{தொடரைக்} கருதலாம்:
\[
1, 1.4, 1.414, 1.4142, 1.41421,\ldots
\]

% SEGMENT OLP-0048-S04
``மேன்மேலும் சிறந்த தோராயங்கள்'' மூலம் மெய்யெண்களைக் கருதலாம்
என்ற கருத்து, மற்றொரு தொடர் உட்பார்வைகளுக்கு அடிப்படையைத் தருகிறது
($\Int$ இலிருந்து $\Rat$ ஐயோ, $\Nat$ இலிருந்து $\Int$ ஐயோ
கட்டமைக்கப் பயன்படுத்திய உணர்தல்களைப் போன்றவை).
அடிப்படை உட்பார்வைகள் இவை:
\begin{enumerate}
	\item ஒவ்வொரு மெய்யெண்ணையும் ஒரு (முடிவுறாததாக இருக்கக்கூடிய)
	தசம விரிவாக்கமாக எழுதலாம்.
	\item ஒரு (முடிவுறாததாக இருக்கக்கூடிய) தசம விரிவாக்கத்தால்
	குறியாக்கப்படும் தகவலை, விகிதமுறு எண்களின் ஒரு தொடராலும்
    அதேபோலக் குறியாக்கலாம்.
	\item விகிதமுறு எண்களின் ஒரு தொடரை $\Nat$ இலிருந்து $\Rat$ க்குச்
    செல்லும் ஒரு சார்பாகக் கருதலாம்; தொடரின் $n$ ஆவது விகிதமுறு
    எண்ணாக $f(n)$ ஐ அமைத்தால் போதும்.
\end{enumerate}
நிச்சயமாக, $\Nat$ இலிருந்து $\Rat$ க்குச் செல்லும் \emph{எந்தச்}
சார்பும் ஒரு மெய்யெண்ணை நமக்குத் தராது.
எடுத்துக்காட்டாக, இந்தச் சார்பைக் கருதுக:
\[
	f(n) =\begin{cases}
			1 & \text{எனில் }n\text{ ஒற்றை}\\
			0 &\text{எனில் }n\text{ இரட்டை}
		\end{cases}
\]
இங்குள்ள அடிப்படைச் சிக்கல், $0,1,0,1,0,1,0,\ldots$ என்ற தொடர்
எந்த மெய்யெண்ணையும் ``நெருங்குவதாகத்'' தெரியவில்லை என்பதாகும்.
எனவே ஏதாவது ஒரு மெய்யெண்ணை நெருங்கும் தொடர்களை மட்டும் கருதுவதை
உறுதிசெய்ய, ஏதாவது ஓர் \emph{எல்லையைக்} கொண்ட தொடர்களுக்கு
நமது கவனத்தை வரம்பிட வேண்டும்.

% SEGMENT OLP-0048-S05
\olref[his][set][limits]{sec} இல் எல்லை என்ற கருத்தை ஏற்கெனவே கண்டுள்ளோம்.
ஆனால் அங்கு பயன்படுத்திய அதே வரையறையை \emph{முற்றிலும்}
பயன்படுத்த முடியாது.
அங்குள்ள ``$(\forall \epsilon>0)$'' என்ற கோவை,
மெய்யெண்கள் மீதான அளவையிடலை மறைமுகமாகக் கொண்டிருந்தது;
மேலும் மெய்யெண்கள் மீதான சார்புகளின் எல்லைகளைக் கருதினோம்.
எனவே அந்த வரையறையைப் பயன்படுத்துவது,
மெய்யெண்களை நமக்குத் தரப்பட்டவையாகப் பயன்படுத்துவதாகும்;
ஆனால் அவற்றையே \emph{கட்டமைப்பது} நமது நோக்கம்.
நல்வாய்ப்பாக, நெருங்கிய தொடர்புடைய ஓர் எல்லைக் கருத்தைப் பயன்படுத்தலாம்.
\begin{defn}\ollabel{def:CauchySequence}
  எந்த நேர்ம $\epsilon \in \Rat$ க்கும் $(\exists \ell \in
  \Nat)(\forall m, n > \ell)|f(m) - f(n)| < \epsilon$ என
  இருக்கும்போது, அப்போதும் அப்போது மட்டுமே,
  $f: \Nat \to \Rat$ என்ற சார்பு ஒரு \emph{கோஷி தொடர்} ஆகும்.
\end{defn}

% SEGMENT OLP-0048-S06
ஓர் எல்லையின் பொதுக் கருத்து முன்புபோலவே உள்ளது:
ஒரு குறிப்பிட்ட துல்லிய அளவு ($\epsilon$ ஆல் அளக்கப்படுவது) வேண்டுமெனில்,
பார்க்க வேண்டிய ஒரு ``பகுதி'' உள்ளது ($\ell$ ஐ விடப் பெரிய எந்த உள்ளீடும்).
நமது $1$, $1.4$, $1.414$, $1.4142$, $1.41421$\ldots என்ற தொடருக்கு
ஓர் எல்லை உள்ளது என்பதை எளிதில் காணலாம்:
$\sqrt{2}$ ஐ $\nicefrac{1}{10^{n}}$ என்ற பிழைக்குள் தோராயப்படுத்த
விரும்பினால், $n$ ஆவது உறுப்புக்குப் பிந்தைய எந்த உறுப்பையும் பார்க்கலாம்.

அப்படியானால், ஒரு மெய்யெண் எந்தவொரு கோஷி தொடராகவும் \emph{உள்ளது}
என்று கூறுவது தெளிவான எண்ணமாக இருக்கும்.
ஆனால் $\Int$, $\Rat$ ஆகியவற்றின் கட்டமைப்புகளைப் போலவே,
இது அளவுக்கு அதிகமாக முறைசாராதது:
கொடுக்கப்பட்ட எந்த மெய்யெண்ணையும் பல வெவ்வேறு கோஷி தொடர்கள் குறிக்கின்றன.
இதைக் காண்பதற்கான எளிய வழி இதோ.
ஒரு கோஷி தொடர் $f$ தரப்பட்டால், $g(0)\neq f(0)$ என்பதைத் தவிர்த்து,
$f$ ஐ மிகச்சரியாக ஒத்த அதே சார்பாக $g$ ஐ வரையறுக்கவும்.
முதல் எண்ணுக்குப் பிறகு இரு தொடர்களும் எல்லா இடங்களிலும் ஒத்துப்போவதால்,
\olref{def:CauchySequence} இல் பயன்படுத்திய பொருளில் அவை ஒரே
எல்லையைக் கொண்டுள்ளன என்றும், எனவே ஒரே மெய்யெண்ணை
``வரையறுப்பவையாகக்'' கருத வேண்டும் என்றும் (இறுதியில்) கூற விரும்புவோம்.
ஆகவே இந்தக் கோஷி தொடர்களை உண்மையில் ஒரே மெய்யெண்ணாகக் கருத வேண்டும்.

% SEGMENT OLP-0048-S07
இதன் விளைவாக, கோஷி தொடர்கள் மீது மீண்டும் ஒரு சமானத் தொடர்பை
வரையறுத்து, மெய்யெண்களைச் சமான வகுப்புகளுடன் அடையாளப்படுத்த வேண்டும்.
முதலில், எல்லையில் $0$ ஐ அணுகும் சார்பு என்ற கருத்து தேவை.
எந்தச் சார்பு $h : \Nat \to \Rat$ க்கும்,
எந்த நேர்ம $\epsilon \in \Rat$ க்கும் $(\exists \ell \in
\Nat)(\forall n > \ell)|h(n)| < \epsilon$ என இருக்கும்போது,
அப்போதும் அப்போது மட்டுமே, \emph{$h$ ஆனது $0$ ஐ அணுகுகிறது} என்க.
\footnote{இதனை \olref[his][set][limits]{sec} இல் உள்ள
$\lim_{x \mathord{\rightarrow}\infty}f(x) = 0$ என்ற வரையறையுடன் ஒப்பிடுக.}
மேலும், $f$, $g$ என்பவை $\Nat \to \Rat$ என்ற சார்புகளாக இருக்கும்போது,
$(f-g)(n) = f(n) - g(n)$ என்க. இப்போது பின்வருமாறு வரையறுக்க:
\[
	f \Realequiv g \text{ என்பது $(f-g)$ $0$ ஐ அணுகினால், அப்போது மட்டுமே}.
\]
$\Realequiv$ ஒரு சமானத் தொடர்பு எனச் சரிபார்க்க வேண்டும்; அது அவ்வாறே.
பின்னர் நாம் விரும்பினால், $\Nat \to \Rat$ என்ற எல்லாக் கோஷி
தொடர்களுக்கும் $\Realequiv$ இன் கீழ் கிடைக்கும் சமான வகுப்புகளாக
மெய்யெண்களை வரையறுக்கலாம்.

% SEGMENT OLP-0048-S08
\begin{prob}
ஒவ்வொரு $n$ க்கும் $f(n) = 0$ என்க.
$g(n) = \frac{1}{(n+1)^2}$ என்க.
இரண்டும் கோஷி தொடர்கள் எனவும், உண்மையில் இரு சார்புகளின் எல்லையும்
$0$ எனவும், ஆகவே $f \sim_\Real g$ எனவும் காட்டுக.
\end{prob}

% SEGMENT OLP-0048-S09
இதனைச் செய்தபின், வழக்கம்போல $f$ ஐ ஓர் !!a{element} ஆகக் கொண்ட
சமான வகுப்பை $\equivrep{f}{\Realequiv}$ என எழுதுவோம்.
ஆனால் எளிதாகப் படிப்பதற்கு, பின்வருவதில் கீழ்க்குறியை விட்டுவிட்டு
$\equivrep{f}{}$ என மட்டும் எழுதுவோம்.
மேலும் ஒவ்வொரு $q \in \Rat$ க்கும்,
$q_{\Real} = \equivrep{c_{q}}{}$ என நிர்ணயிக்கிறோம்;
இங்கு $c_{q}$ என்பது அனைத்து $n \in \Nat$ க்கும்
$c_q(n) = q$ என அமையும் மாறிலிச் சார்பு.
பின்னர் மெய்யெண்களின் மீதான அடிப்படைத் தொடர்புகளையும் செயல்களையும்,
எடுத்துக்காட்டாகப் பின்வருமாறு வரையறுக்கிறோம்:
\begin{align*}
	\equivrep{f}{} + \equivrep{g}{} &= 	\equivrep{(f + g)}{} \\
	\equivrep{f}{} \times 	\equivrep{g}{} &= \equivrep{(f \times g)}{}
\end{align*}
இங்கு $(f + g)(n) = f(n) + g(n)$, $(f \times g)(n) = f(n) \times
g(n)$.
நிச்சயமாக, $f$, $g$ ஆகியவை கோஷி தொடர்களாக இருக்கும்போது,
$(f + g)$, $(f-g)$, $(f\times g)$ ஒவ்வொன்றும் கோஷி தொடராக உள்ளதா
என்றும் சரிபார்க்க வேண்டும்; அவை அவ்வாறே உள்ளன,
இதனை உங்களுக்கே விட்டுவிடுகிறோம்.

% SEGMENT OLP-0048-S10
இறுதியாக, வரிசை என்ற ஒரு கருத்தை வரையறுக்கிறோம்.
$\equivrep{f}{}\neq 0_\Real$, $(\exists
\ell \in \Nat)(\forall n > \ell)0 < f(n)$ ஆகிய இரண்டும்
நிறைவேறும்போது, அப்போதும் அப்போது மட்டுமே,
$\equivrep{f}{}$ \emph{நேர்மமானது} என்க.
பின்னர் $\equivrep{(g - f)}{}$ நேர்மமாக இருக்கும்போது,
அப்போதும் அப்போது மட்டுமே,
$\equivrep{f}{} < \equivrep{g}{}$ என்க.
இது நன்கு வரையறுக்கப்பட்டுள்ளது எனச் சரிபார்க்க வேண்டும்
(அதாவது சமான வகுப்பிலிருந்து தேர்ந்தெடுக்கும் ``பிரதிநிதிச்'' சார்பை
இது சார்ந்திருக்கவில்லை எனச் சரிபார்க்க வேண்டும்).
%\begin{align*}
%	\equivrep{f}{} \leq \equivrep{g}{} \text{ ஆக இருக்கும்போது, அப்போதும் அப்போது மட்டுமே, }&\equivrep{f}{} = \equivrep{g}{} \text{ அல்லது }\\
%	&(\exists \ell \in \Nat)(\forall n \in \Nat)(\ell < n \lif f(n) \leq g(n))
%\end{align*}
ஆனால் இதைச் செய்தபின், இவை சரியான இயற்கணிதப் பண்புகளைத் தருகின்றன
எனக் காட்டுவது மிகவும் எளிது; அதாவது:

% SEGMENT OLP-0048-S11
\begin{thm}\ollabel{thm:cauchyorderedfield}
கோஷி தொடர்களின் சமான வகுப்புகள் ஒரு வரிசைப்படுத்தப்பட்ட புலத்தை அமைக்கின்றன.
\end{thm}

% SEGMENT OLP-0048-S12
\begin{proof}
பயிற்சி.
\end{proof}

% SEGMENT OLP-0048-S13
\begin{prob}
கோஷி தொடர்களின் சமான வகுப்புகள் ஒரு வரிசைப்படுத்தப்பட்ட புலத்தை
அமைக்கின்றன என நிறுவுக.
\end{prob}

% SEGMENT OLP-0048-S14
இவ்வாறு கட்டமைக்கப்பட்ட மெய்யெண்களுக்கு முழுமைப் பண்பு உள்ளது என
நிறுவுவது மேலும் கடினம்; எனவே நிறுவலைத் தருவோம்.

% SEGMENT OLP-0048-S15
\begin{thm}
மேல்வரம்பைக் கொண்ட கோஷி தொடர்களின் ஒவ்வொரு வெற்றற்ற கணமும்
ஒரு மீச்சிறு மேல்வரம்பைக் கொண்டுள்ளது.
\end{thm}

% SEGMENT OLP-0048-S16
\begin{proof}[நிறுவல் வரைவு]
$S$ என்பது மேல்வரம்பைக் கொண்ட கோஷி தொடர்களின் எந்தவொரு வெற்றற்ற
கணமாகவும் இருக்கட்டும்.
அப்போது $p_{\Real}$ என்பது $S$ இன் மேல்வரம்பாக அமையும் ஏதோ ஒரு
$p \in \Rat$ உள்ளது.
$r \in S$ என்க; அப்போது $q_{\Real} < r$ என அமையும் ஏதோ ஒரு
$q \in \Rat$ உள்ளது.
எனவே $S$ க்கு ஒரு மீச்சிறு மேல்வரம்பு இருந்தால்,
அது $q_\Real$, $p_\Real$ ஆகியவற்றுக்கு இடையில் (அவை உட்பட) இருக்கும்.

மீச்சிறு மேல்வரம்பை ஒரே நேரத்தில் கீழிருந்தும் மேலிருந்தும் அணுகி,
அதனை நெருங்குவோம்.
குறிப்பாக, $f$ மேலிருந்தும் $g$ கீழிருந்தும் மீச்சிறு மேல்வரம்பை
நெருங்க வேண்டும் என்ற நோக்கில், $f, g \colon
\Nat \to \Rat$ என்ற இரண்டு சார்புகளை வரையறுக்கிறோம்.
பின்வருமாறு வரையறுத்துத் தொடங்குகிறோம்:
\begin{align*}
	f(0) &= p \\
	g(0) &= q
\end{align*}
பின்னர் $a_n = \frac{f(n) + g(n)}{2}$ ஆக இருக்கும்போது,
பின்வருமாறு அமைக்க:
\footnote{இது ஒரு மறுநிகழ்வு வரையறை.
ஆனால் மறுநிகழ்வு வரையறைகள் ஏற்கத்தக்கவை எனக் கருதுவதற்கு
\emph{இன்னும்} எந்தக் காரணமும் தரவில்லை.}
\begin{align*}
	f(n+1) &=
	\begin{cases}
		a_n &\text{எனில் }(\forall h \in S)\equivrep{h}{} \leq (a_n)_\Real\\
		f(n)&\text{மற்ற நிலைகளில்}
	\end{cases}\\
	g(n+1) &=
	\begin{cases}
		a_n &\text{எனில் }(\exists h \in S)\equivrep{h}{} \geq (a_n)_\Real\\
	 	g(n) &\text{மற்ற நிலைகளில்}
	\end{cases}
\end{align*}
$f$, $g$ இரண்டும் கோஷி தொடர்கள்.
(இதனை எளிதாகச் சரிபார்க்கலாம்; ஆனால் ஒரு பயிற்சியாக விடுகிறோம்.)
ஒவ்வொரு படியிலும் $f$, $g$ ஆகியவற்றுக்கு இடையிலான வேறுபாடு
பாதியாகக் குறைவதால், $(f-g)$ என்ற சார்பு $0$ ஐ அணுகுகிறது.
எனவே $\equivrep{f}{} = \equivrep{g}{}$.

முதலில் $\equivrep{f}{}$ என்பது $S$ இன் ஒரு மேல்வரம்பு எனக் காட்டுகிறோம்;
அதாவது $(\forall h \in S)\equivrep{h}{} \leq \equivrep{f}{}$.
(தொடரும்போது \olref{thm:cauchyorderedfield} ஐப் பயன்படுத்துவோம்.)
$h \in S$ என்க; முரண் பெறும் நோக்கில்,
$\equivrep{f}{} < \equivrep{h}{}$ எனக் கொள்க;
எனவே $0_\Real < \equivrep{(h-f)}{}$.
$f$ என்பது ஒருபோக்காகக் குறையும் கோஷி தொடர் என்பதால்,
$\equivrep{(c_{f(n)} - f)}{} < \equivrep{(h-f)}{}$ என அமையும்
ஏதோ ஒரு $n \in \Nat$ உள்ளது. எனவே:
\[
	(f(n))_\Real = \equivrep{c_{f(n)}}{} < \equivrep{f}{} + \equivrep{(h-f)}{} = \equivrep{h}{},
\]
ஆனால் கட்டுமானத்தின்படி $\equivrep{h}{} \leq (f(n))_\Real$;
இதற்கு அது முரணாகிறது.

அடுத்து $\equivrep{f}{} = \equivrep{g}{}$ என்பது $S$ இன்
\emph{மீச்சிறு} மேல்வரம்பு எனக் காட்டுகிறோம்.
$j$ என்பது எந்தவொரு கோஷி தொடராகவும் இருக்கட்டும்;
$\equivrep{j}{} < \equivrep{g}{}$ எனக் கொள்க.
மேலே செய்ததுபோலக் காரணப்படுத்தினால் ($g$ \emph{அதிகரிக்கிறது}
என்ற உண்மையைப் பயன்படுத்தி),
$\equivrep{j}{} < (g(n))_\Real$ என அமையும் $n \in \Nat$ உள்ளது.
ஆனால் கட்டுமானத்தின்படி $(g(n))_\Real \leq \equivrep{h}{}$ என
அமையும் $h \in S$ உள்ளது; எனவே $\equivrep{j}{} < \equivrep{h}{}$;
ஆகவே $\equivrep{j}{}$ என்பது $S$ இன் மேல்வரம்பு அல்ல.
\end{proof}

\end{document}
